% Options for packages loaded elsewhere
\PassOptionsToPackage{unicode}{hyperref}
\PassOptionsToPackage{hyphens}{url}
%
\documentclass[
  ignorenonframetext,
]{beamer}
\usepackage{pgfpages}
\setbeamertemplate{caption}[numbered]
\setbeamertemplate{caption label separator}{: }
\setbeamercolor{caption name}{fg=normal text.fg}
\beamertemplatenavigationsymbolsempty
% Prevent slide breaks in the middle of a paragraph
\widowpenalties 1 10000
\raggedbottom
\setbeamertemplate{part page}{
  \centering
  \begin{beamercolorbox}[sep=16pt,center]{part title}
    \usebeamerfont{part title}\insertpart\par
  \end{beamercolorbox}
}
\setbeamertemplate{section page}{
  \centering
  \begin{beamercolorbox}[sep=12pt,center]{part title}
    \usebeamerfont{section title}\insertsection\par
  \end{beamercolorbox}
}
\setbeamertemplate{subsection page}{
  \centering
  \begin{beamercolorbox}[sep=8pt,center]{part title}
    \usebeamerfont{subsection title}\insertsubsection\par
  \end{beamercolorbox}
}
\AtBeginPart{
  \frame{\partpage}
}
\AtBeginSection{
  \ifbibliography
  \else
    \frame{\sectionpage}
  \fi
}
\AtBeginSubsection{
  \frame{\subsectionpage}
}
\usepackage{amsmath,amssymb}
\usepackage{lmodern}
\usepackage{iftex}
\ifPDFTeX
  \usepackage[T1]{fontenc}
  \usepackage[utf8]{inputenc}
  \usepackage{textcomp} % provide euro and other symbols
\else % if luatex or xetex
  \usepackage{unicode-math}
  \defaultfontfeatures{Scale=MatchLowercase}
  \defaultfontfeatures[\rmfamily]{Ligatures=TeX,Scale=1}
\fi
\usetheme[]{Montpellier}
% Use upquote if available, for straight quotes in verbatim environments
\IfFileExists{upquote.sty}{\usepackage{upquote}}{}
\IfFileExists{microtype.sty}{% use microtype if available
  \usepackage[]{microtype}
  \UseMicrotypeSet[protrusion]{basicmath} % disable protrusion for tt fonts
}{}
\makeatletter
\@ifundefined{KOMAClassName}{% if non-KOMA class
  \IfFileExists{parskip.sty}{%
    \usepackage{parskip}
  }{% else
    \setlength{\parindent}{0pt}
    \setlength{\parskip}{6pt plus 2pt minus 1pt}}
}{% if KOMA class
  \KOMAoptions{parskip=half}}
\makeatother
\usepackage{xcolor}
\IfFileExists{xurl.sty}{\usepackage{xurl}}{} % add URL line breaks if available
\IfFileExists{bookmark.sty}{\usepackage{bookmark}}{\usepackage{hyperref}}
\hypersetup{
  hidelinks,
  pdfcreator={LaTeX via pandoc}}
\urlstyle{same} % disable monospaced font for URLs
\newif\ifbibliography
\usepackage{color}
\usepackage{fancyvrb}
\newcommand{\VerbBar}{|}
\newcommand{\VERB}{\Verb[commandchars=\\\{\}]}
\DefineVerbatimEnvironment{Highlighting}{Verbatim}{commandchars=\\\{\}}
% Add ',fontsize=\small' for more characters per line
\usepackage{framed}
\definecolor{shadecolor}{RGB}{248,248,248}
\newenvironment{Shaded}{\begin{snugshade}}{\end{snugshade}}
\newcommand{\AlertTok}[1]{\textcolor[rgb]{0.94,0.16,0.16}{#1}}
\newcommand{\AnnotationTok}[1]{\textcolor[rgb]{0.56,0.35,0.01}{\textbf{\textit{#1}}}}
\newcommand{\AttributeTok}[1]{\textcolor[rgb]{0.77,0.63,0.00}{#1}}
\newcommand{\BaseNTok}[1]{\textcolor[rgb]{0.00,0.00,0.81}{#1}}
\newcommand{\BuiltInTok}[1]{#1}
\newcommand{\CharTok}[1]{\textcolor[rgb]{0.31,0.60,0.02}{#1}}
\newcommand{\CommentTok}[1]{\textcolor[rgb]{0.56,0.35,0.01}{\textit{#1}}}
\newcommand{\CommentVarTok}[1]{\textcolor[rgb]{0.56,0.35,0.01}{\textbf{\textit{#1}}}}
\newcommand{\ConstantTok}[1]{\textcolor[rgb]{0.00,0.00,0.00}{#1}}
\newcommand{\ControlFlowTok}[1]{\textcolor[rgb]{0.13,0.29,0.53}{\textbf{#1}}}
\newcommand{\DataTypeTok}[1]{\textcolor[rgb]{0.13,0.29,0.53}{#1}}
\newcommand{\DecValTok}[1]{\textcolor[rgb]{0.00,0.00,0.81}{#1}}
\newcommand{\DocumentationTok}[1]{\textcolor[rgb]{0.56,0.35,0.01}{\textbf{\textit{#1}}}}
\newcommand{\ErrorTok}[1]{\textcolor[rgb]{0.64,0.00,0.00}{\textbf{#1}}}
\newcommand{\ExtensionTok}[1]{#1}
\newcommand{\FloatTok}[1]{\textcolor[rgb]{0.00,0.00,0.81}{#1}}
\newcommand{\FunctionTok}[1]{\textcolor[rgb]{0.00,0.00,0.00}{#1}}
\newcommand{\ImportTok}[1]{#1}
\newcommand{\InformationTok}[1]{\textcolor[rgb]{0.56,0.35,0.01}{\textbf{\textit{#1}}}}
\newcommand{\KeywordTok}[1]{\textcolor[rgb]{0.13,0.29,0.53}{\textbf{#1}}}
\newcommand{\NormalTok}[1]{#1}
\newcommand{\OperatorTok}[1]{\textcolor[rgb]{0.81,0.36,0.00}{\textbf{#1}}}
\newcommand{\OtherTok}[1]{\textcolor[rgb]{0.56,0.35,0.01}{#1}}
\newcommand{\PreprocessorTok}[1]{\textcolor[rgb]{0.56,0.35,0.01}{\textit{#1}}}
\newcommand{\RegionMarkerTok}[1]{#1}
\newcommand{\SpecialCharTok}[1]{\textcolor[rgb]{0.00,0.00,0.00}{#1}}
\newcommand{\SpecialStringTok}[1]{\textcolor[rgb]{0.31,0.60,0.02}{#1}}
\newcommand{\StringTok}[1]{\textcolor[rgb]{0.31,0.60,0.02}{#1}}
\newcommand{\VariableTok}[1]{\textcolor[rgb]{0.00,0.00,0.00}{#1}}
\newcommand{\VerbatimStringTok}[1]{\textcolor[rgb]{0.31,0.60,0.02}{#1}}
\newcommand{\WarningTok}[1]{\textcolor[rgb]{0.56,0.35,0.01}{\textbf{\textit{#1}}}}
\setlength{\emergencystretch}{3em} % prevent overfull lines
\providecommand{\tightlist}{%
  \setlength{\itemsep}{0pt}\setlength{\parskip}{0pt}}
\setcounter{secnumdepth}{-\maxdimen} % remove section numbering
%\documentclass[unicode,10pt]{beamer}% 'unicode'が必要
\usepackage{luatexja}% 日本語を使う場合は必須
%\usepackage[japanese]{babel}	
%\usefonttheme[onlymath]{serif}
%\usepackage[T1]{fontenc} %これを使うと、ギリシャ文字が出ない
%\usepackage{textcomp}
%\usepackage[scale=1.0]{tgheros} % Sans serif font
%\usepackage[scaled]{beramono} % Monospace font
%\usepackage{luatexja-otf}
%\renewcommand{\kanjifamilydefault}{\gtdefault}
%\usepackage[haranoaji]{luatexja-preset}

% スライド番号の表示 %%%%%%%%%%%%%%%%%%%
\makeatletter
\setbeamertemplate{footline}{%
  \leavevmode%
  \hbox{%
  \begin{beamercolorbox}[wd=.5\paperwidth,ht=2.25ex,dp=1ex,right]{date in head/foot}%
    \insertframenumber{} \hspace*{2ex} % new version without total frames
  \end{beamercolorbox}}%
  \vskip0pt%
}
\makeatother
% スライド番号の表示 %%%%%%%%%%%%%%%%%%%
\ifLuaTeX
  \usepackage{selnolig}  % disable illegal ligatures
\fi

\title{第4章: 予測\\
(4.1. 選挙結果の予測)}
\subtitle{今井耕介 著\\
『社会科学のためのデータ分析入門 (QSS)』}
\author{}
\date{\vspace{-2.5em}2026-03-09}

\begin{document}
\frame{\titlepage}

\hypertarget{ux9078ux6319ux7d50ux679cux306eux4e88ux6e2c}{%
\section{4.1
選挙結果の予測}\label{ux9078ux6319ux7d50ux679cux306eux4e88ux6e2c}}

\begin{frame}{予測とは}
\protect\hypertarget{ux4e88ux6e2cux3068ux306f}{}
\begin{itemize}
\tightlist
\item
  \textbf{予測}:
  既存のデータ（訓練データ）からパターンを学び、未知のデータや未来の結果を推測すること。
\item
  社会科学の例：

  \begin{itemize}
  \tightlist
  \item
    世論調査から選挙結果を予測する。
  \item
    経済指標から景気後退を予測する。
  \end{itemize}
\item
  本節では、Rのプログラミングの基礎（ループと条件分岐）を学びながら、2008年米大統領選の予測に取り組みます。
\end{itemize}
\end{frame}

\hypertarget{rux3067ux306eux30ebux30fcux30d7-loops}{%
\section{4.1.1 Rでのループ
(Loops)}\label{rux3067ux306eux30ebux30fcux30d7-loops}}

\begin{frame}[fragile]{for ループの基本: 準備}
\protect\hypertarget{for-ux30ebux30fcux30d7ux306eux57faux672c-ux6e96ux5099}{}
\begin{itemize}
\tightlist
\item
  同じ操作を繰り返し実行したい時に使用します。
\end{itemize}

\scriptsize

\begin{Shaded}
\begin{Highlighting}[]
\NormalTok{values }\OtherTok{\textless{}{-}} \FunctionTok{c}\NormalTok{(}\DecValTok{2}\NormalTok{, }\DecValTok{4}\NormalTok{, }\DecValTok{6}\NormalTok{) }\CommentTok{\# 計算対象のベクトルを作成}
\NormalTok{n }\OtherTok{\textless{}{-}} \FunctionTok{length}\NormalTok{(values)   }\CommentTok{\# ベクトルの長さを取得}
\NormalTok{results }\OtherTok{\textless{}{-}} \FunctionTok{rep}\NormalTok{(}\ConstantTok{NA}\NormalTok{, n) }\CommentTok{\# 結果保存用の空ベクトルをNAで初期化}
\end{Highlighting}
\end{Shaded}

\normalsize
\end{frame}

\begin{frame}[fragile]{for ループの基本: 実行}
\protect\hypertarget{for-ux30ebux30fcux30d7ux306eux57faux672c-ux5b9fux884c}{}
\scriptsize

\begin{Shaded}
\begin{Highlighting}[]
\CommentTok{\# i が 1 から n まで順に変わるループ}
\ControlFlowTok{for}\NormalTok{ (i }\ControlFlowTok{in} \DecValTok{1}\SpecialCharTok{:}\NormalTok{n) \{}
\NormalTok{    results[i] }\OtherTok{\textless{}{-}}\NormalTok{ values[i] }\SpecialCharTok{*} \DecValTok{2} \CommentTok{\# valuesのi番目を2倍してresultsのi番目に代入}
    \FunctionTok{cat}\NormalTok{(values[i], }\StringTok{"times 2 is equal to"}\NormalTok{, results[i], }\StringTok{"}\SpecialCharTok{\textbackslash{}n}\StringTok{"}\NormalTok{) }\CommentTok{\# 計算過程を表示}
\NormalTok{\}}
\end{Highlighting}
\end{Shaded}

\begin{verbatim}
## 2 times 2 is equal to 4 
## 4 times 2 is equal to 8 
## 6 times 2 is equal to 12
\end{verbatim}

\begin{Shaded}
\begin{Highlighting}[]
\NormalTok{results }\CommentTok{\# 最終的な結果を表示}
\end{Highlighting}
\end{Shaded}

\begin{verbatim}
## [1]  4  8 12
\end{verbatim}

\normalsize
\end{frame}

\hypertarget{rux3067ux306eux6761ux4ef6ux5206ux5c90-conditional-statements}{%
\section{4.1.2 Rでの条件分岐 (Conditional
Statements)}\label{rux3067ux306eux6761ux4ef6ux5206ux5c90-conditional-statements}}

\begin{frame}[fragile]{if 文の基本: 準備}
\protect\hypertarget{if-ux6587ux306eux57faux672c-ux6e96ux5099}{}
\begin{itemize}
\tightlist
\item
  条件によって実行するコードを切り替えます。
\end{itemize}

\scriptsize

\begin{Shaded}
\begin{Highlighting}[]
\NormalTok{values }\OtherTok{\textless{}{-}} \DecValTok{1}\SpecialCharTok{:}\DecValTok{5}                     \CommentTok{\# 1から5までの連続数を作成}
\NormalTok{results }\OtherTok{\textless{}{-}} \FunctionTok{rep}\NormalTok{(}\ConstantTok{NA}\NormalTok{, }\FunctionTok{length}\NormalTok{(values)) }\CommentTok{\# 結果保存用の空ベクトルを作成}
\end{Highlighting}
\end{Shaded}

\normalsize
\end{frame}

\begin{frame}[fragile]{if 文の基本: 実行}
\protect\hypertarget{if-ux6587ux306eux57faux672c-ux5b9fux884c}{}
\scriptsize

\begin{Shaded}
\begin{Highlighting}[]
\ControlFlowTok{for}\NormalTok{ (i }\ControlFlowTok{in} \DecValTok{1}\SpecialCharTok{:}\FunctionTok{length}\NormalTok{(values)) \{}
\NormalTok{    x }\OtherTok{\textless{}{-}}\NormalTok{ values[i] }\CommentTok{\# 現在の要素を取り出す}
    \ControlFlowTok{if}\NormalTok{ (x }\SpecialCharTok{\%\%} \DecValTok{2} \SpecialCharTok{==} \DecValTok{0}\NormalTok{) \{ }\CommentTok{\# 2で割った余りが0（偶数）の場合}
        \FunctionTok{cat}\NormalTok{(x, }\StringTok{"is even. Adding itself.}\SpecialCharTok{\textbackslash{}n}\StringTok{"}\NormalTok{) }\CommentTok{\# 偶数であることを表示}
\NormalTok{        results[i] }\OtherTok{\textless{}{-}}\NormalTok{ x }\SpecialCharTok{+}\NormalTok{ x                 }\CommentTok{\# 自分自身を足す}
\NormalTok{    \} }\ControlFlowTok{else}\NormalTok{ \{           }\CommentTok{\# 奇数の場合}
        \FunctionTok{cat}\NormalTok{(x, }\StringTok{"is odd. Multiplying itself.}\SpecialCharTok{\textbackslash{}n}\StringTok{"}\NormalTok{) }\CommentTok{\# 奇数であることを表示}
\NormalTok{        results[i] }\OtherTok{\textless{}{-}}\NormalTok{ x }\SpecialCharTok{*}\NormalTok{ x                    }\CommentTok{\# 自分自身を掛ける}
\NormalTok{    \}}
\NormalTok{\}}
\end{Highlighting}
\end{Shaded}

\begin{verbatim}
## 1 is odd. Multiplying itself.
## 2 is even. Adding itself.
## 3 is odd. Multiplying itself.
## 4 is even. Adding itself.
## 5 is odd. Multiplying itself.
\end{verbatim}

\begin{Shaded}
\begin{Highlighting}[]
\NormalTok{results }\CommentTok{\# 結果を表示}
\end{Highlighting}
\end{Shaded}

\begin{verbatim}
## [1]  1  4  9  8 25
\end{verbatim}

\normalsize
\end{frame}

\hypertarget{ux4e16ux8ad6ux8abfux67fbux306bux3088ux308bux4e88ux6e2c}{%
\section{4.1.3
世論調査による予測}\label{ux4e16ux8ad6ux8abfux67fbux306bux3088ux308bux4e88ux6e2c}}

\begin{frame}[fragile]{選挙データの読み込み (1)}
\protect\hypertarget{ux9078ux6319ux30c7ux30fcux30bfux306eux8aadux307fux8fbcux307f-1}{}
\begin{itemize}
\tightlist
\item
  2008年大統領選の確定結果 (\texttt{pres08.csv}) と世論調査
  (\texttt{polls08.csv}) を読み込みます。
\end{itemize}

\scriptsize

\begin{Shaded}
\begin{Highlighting}[]
\NormalTok{pres08 }\OtherTok{\textless{}{-}} \FunctionTok{read.csv}\NormalTok{(}\StringTok{"pres08.csv"}\NormalTok{)   }\CommentTok{\# 選挙結果データを読み込む}
\NormalTok{polls08 }\OtherTok{\textless{}{-}} \FunctionTok{read.csv}\NormalTok{(}\StringTok{"polls08.csv"}\NormalTok{) }\CommentTok{\# 世論調査データを読み込む}

\CommentTok{\# （参考）URLから直接読み込むことも可能}
\CommentTok{\# pres08 \textless{}{-} read.csv("https://ayumu{-}tanaka.github.io/QSS/QSS\_Data/pres08.csv")}
\CommentTok{\# polls08 \textless{}{-} read.csv("https://ayumu{-}tanaka.github.io/QSS/QSS\_Data/polls08.csv")}
\end{Highlighting}
\end{Shaded}

\normalsize
\end{frame}

\begin{frame}[fragile]{データの準備 (2)}
\protect\hypertarget{ux30c7ux30fcux30bfux306eux6e96ux5099-2}{}
\begin{itemize}
\tightlist
\item
  得票率のマージン計算と、日付データの処理を行います。
\end{itemize}

\scriptsize

\begin{Shaded}
\begin{Highlighting}[]
\NormalTok{polls08}\SpecialCharTok{$}\NormalTok{margin }\OtherTok{\textless{}{-}}\NormalTok{ polls08}\SpecialCharTok{$}\NormalTok{Obama }\SpecialCharTok{{-}}\NormalTok{ polls08}\SpecialCharTok{$}\NormalTok{McCain }\CommentTok{\# 世論調査の支持率差を計算}
\NormalTok{pres08}\SpecialCharTok{$}\NormalTok{margin }\OtherTok{\textless{}{-}}\NormalTok{ pres08}\SpecialCharTok{$}\NormalTok{Obama }\SpecialCharTok{{-}}\NormalTok{ pres08}\SpecialCharTok{$}\NormalTok{McCain   }\CommentTok{\# 実際の選挙結果の得票率差を計算}
\NormalTok{polls08}\SpecialCharTok{$}\NormalTok{middate }\OtherTok{\textless{}{-}} \FunctionTok{as.Date}\NormalTok{(polls08}\SpecialCharTok{$}\NormalTok{middate)       }\CommentTok{\# 文字列を日付型に変換}
\CommentTok{\# 選挙当日（2008{-}11{-}04）までの日数を計算}
\NormalTok{polls08}\SpecialCharTok{$}\NormalTok{DaysToElection }\OtherTok{\textless{}{-}} \FunctionTok{as.Date}\NormalTok{(}\StringTok{"2008{-}11{-}04"}\NormalTok{) }\SpecialCharTok{{-}}\NormalTok{ polls08}\SpecialCharTok{$}\NormalTok{middate}
\end{Highlighting}
\end{Shaded}

\normalsize
\end{frame}

\begin{frame}[fragile]{州ごとの最新調査に基づく予測: 準備}
\protect\hypertarget{ux5ddeux3054ux3068ux306eux6700ux65b0ux8abfux67fbux306bux57faux3065ux304fux4e88ux6e2c-ux6e96ux5099}{}
\begin{itemize}
\tightlist
\item
  50州＋DCごとに、選挙に最も近い日の調査結果を平均して予測値とします。
\end{itemize}

\scriptsize

\begin{Shaded}
\begin{Highlighting}[]
\NormalTok{st.names }\OtherTok{\textless{}{-}} \FunctionTok{unique}\NormalTok{(polls08}\SpecialCharTok{$}\NormalTok{state)    }\CommentTok{\# 州名のリスト（重複なし）を作成}
\NormalTok{poll.pred }\OtherTok{\textless{}{-}} \FunctionTok{rep}\NormalTok{(}\ConstantTok{NA}\NormalTok{, }\DecValTok{51}\NormalTok{)             }\CommentTok{\# 51州分の予測値を保存するベクトル}
\FunctionTok{names}\NormalTok{(poll.pred) }\OtherTok{\textless{}{-}} \FunctionTok{as.character}\NormalTok{(st.names) }\CommentTok{\# ベクトルの各要素に州名をつける}
\end{Highlighting}
\end{Shaded}

\normalsize
\end{frame}

\begin{frame}[fragile]{州ごとの最新調査に基づく予測: ループ処理}
\protect\hypertarget{ux5ddeux3054ux3068ux306eux6700ux65b0ux8abfux67fbux306bux57faux3065ux304fux4e88ux6e2c-ux30ebux30fcux30d7ux51e6ux7406}{}
\scriptsize

\begin{Shaded}
\begin{Highlighting}[]
\ControlFlowTok{for}\NormalTok{ (i }\ControlFlowTok{in} \DecValTok{1}\SpecialCharTok{:}\DecValTok{51}\NormalTok{)\{}
    \CommentTok{\# i番目の州のデータのみを抽出}
\NormalTok{    state.data }\OtherTok{\textless{}{-}} \FunctionTok{subset}\NormalTok{(polls08, }\AttributeTok{subset =}\NormalTok{ (state }\SpecialCharTok{==}\NormalTok{ st.names[i]))}
    \CommentTok{\# その州の中で選挙日に最も近い（DaysToElectionが最小）データを抽出}
\NormalTok{    latest }\OtherTok{\textless{}{-}} \FunctionTok{subset}\NormalTok{(state.data, DaysToElection }\SpecialCharTok{==} \FunctionTok{min}\NormalTok{(DaysToElection))}
    \CommentTok{\# 最新調査の支持率差の平均を予測値として保存}
\NormalTok{    poll.pred[i] }\OtherTok{\textless{}{-}} \FunctionTok{mean}\NormalTok{(latest}\SpecialCharTok{$}\NormalTok{margin)}
\NormalTok{\}}
\FunctionTok{head}\NormalTok{(poll.pred) }\CommentTok{\# 予測値の先頭部分を表示}
\end{Highlighting}
\end{Shaded}

\begin{verbatim}
##    AL    AK    AZ    AR    CA    CO 
## -25.0 -19.0  -2.5  -7.0  24.0   7.0
\end{verbatim}

\normalsize
\end{frame}

\begin{frame}[fragile]{予測誤差の分析}
\protect\hypertarget{ux4e88ux6e2cux8aa4ux5deeux306eux5206ux6790}{}
\begin{itemize}
\tightlist
\item
  実際の選挙結果と予測値の差を確認します。
\end{itemize}

\scriptsize

\begin{Shaded}
\begin{Highlighting}[]
\NormalTok{errors }\OtherTok{\textless{}{-}}\NormalTok{ pres08}\SpecialCharTok{$}\NormalTok{margin }\SpecialCharTok{{-}}\NormalTok{ poll.pred }\CommentTok{\# 実際の得票率差と予測値の差（誤差）を計算}
\FunctionTok{names}\NormalTok{(errors) }\OtherTok{\textless{}{-}}\NormalTok{ st.names           }\CommentTok{\# 誤差ベクトルに州名を付与}
\FunctionTok{mean}\NormalTok{(errors)                        }\CommentTok{\# 平均予測誤差を表示}
\end{Highlighting}
\end{Shaded}

\begin{verbatim}
## [1] 1.062092
\end{verbatim}

\begin{Shaded}
\begin{Highlighting}[]
\FunctionTok{sqrt}\NormalTok{(}\FunctionTok{mean}\NormalTok{(errors}\SpecialCharTok{\^{}}\DecValTok{2}\NormalTok{))                }\CommentTok{\# RMSE（二乗平均平方根誤差）を表示}
\end{Highlighting}
\end{Shaded}

\begin{verbatim}
## [1] 5.90894
\end{verbatim}

\normalsize
\end{frame}

\begin{frame}[fragile]{予測誤差の分布: コード}
\protect\hypertarget{ux4e88ux6e2cux8aa4ux5deeux306eux5206ux5e03-ux30b3ux30fcux30c9}{}
\scriptsize

\begin{Shaded}
\begin{Highlighting}[]
\CommentTok{\# 予測誤差のヒストグラムを描画}
\FunctionTok{hist}\NormalTok{(errors, }\AttributeTok{freq =} \ConstantTok{FALSE}\NormalTok{, }\AttributeTok{ylim =} \FunctionTok{c}\NormalTok{(}\DecValTok{0}\NormalTok{, }\FloatTok{0.08}\NormalTok{),}
     \AttributeTok{main =} \StringTok{"Poll Prediction Error"}\NormalTok{,}
     \AttributeTok{xlab =} \StringTok{"Error in Predicted Margin (percentage points)"}\NormalTok{)}

\CommentTok{\# 平均誤差の位置に垂直な赤い点線を追加}
\FunctionTok{abline}\NormalTok{(}\AttributeTok{v =} \FunctionTok{mean}\NormalTok{(errors), }\AttributeTok{lty =} \StringTok{"dashed"}\NormalTok{, }\AttributeTok{col =} \StringTok{"red"}\NormalTok{)}
\end{Highlighting}
\end{Shaded}

\normalsize
\end{frame}

\begin{frame}{予測誤差の分布: 描画結果}
\protect\hypertarget{ux4e88ux6e2cux8aa4ux5deeux306eux5206ux5e03-ux63cfux753bux7d50ux679c}{}
\scriptsize

\begin{center}\includegraphics[height=0.75\textheight]{QSS_Ch04_1_files/figure-beamer/unnamed-chunk-11-1} \end{center}
\normalsize
\end{frame}

\begin{frame}[fragile]{予測と実際の結果の比較: コード}
\protect\hypertarget{ux4e88ux6e2cux3068ux5b9fux969bux306eux7d50ux679cux306eux6bd4ux8f03-ux30b3ux30fcux30c9}{}
\scriptsize

\begin{Shaded}
\begin{Highlighting}[]
\CommentTok{\# 空のプロット軸を作成}
\FunctionTok{plot}\NormalTok{(poll.pred, pres08}\SpecialCharTok{$}\NormalTok{margin, }\AttributeTok{type =} \StringTok{"n"}\NormalTok{,}
     \AttributeTok{xlim =} \FunctionTok{c}\NormalTok{(}\SpecialCharTok{{-}}\DecValTok{40}\NormalTok{, }\DecValTok{90}\NormalTok{), }\AttributeTok{ylim =} \FunctionTok{c}\NormalTok{(}\SpecialCharTok{{-}}\DecValTok{40}\NormalTok{, }\DecValTok{90}\NormalTok{), }
     \AttributeTok{xlab =} \StringTok{"Predicted Margin (Polls)"}\NormalTok{, }\AttributeTok{ylab =} \StringTok{"Actual Margin (Election)"}\NormalTok{)}

\CommentTok{\# 各州の値を略称（州名）でプロット}
\FunctionTok{text}\NormalTok{(}\AttributeTok{x =}\NormalTok{ poll.pred, }\AttributeTok{y =}\NormalTok{ pres08}\SpecialCharTok{$}\NormalTok{margin, }\AttributeTok{labels =}\NormalTok{ pres08}\SpecialCharTok{$}\NormalTok{state, }\AttributeTok{col =} \StringTok{"blue"}\NormalTok{)}

\CommentTok{\# 予測と実績が一致する45度線を追加}
\FunctionTok{abline}\NormalTok{(}\AttributeTok{a =} \DecValTok{0}\NormalTok{, }\AttributeTok{b =} \DecValTok{1}\NormalTok{, }\AttributeTok{lty =} \StringTok{"dashed"}\NormalTok{)}
\FunctionTok{abline}\NormalTok{(}\AttributeTok{v =} \DecValTok{0}\NormalTok{) }\CommentTok{\# 垂直線（0）を追加}
\FunctionTok{abline}\NormalTok{(}\AttributeTok{h =} \DecValTok{0}\NormalTok{) }\CommentTok{\# 水平線（0）を追加}
\end{Highlighting}
\end{Shaded}

\normalsize
\end{frame}

\begin{frame}{予測と実際の結果の比較: 描画結果}
\protect\hypertarget{ux4e88ux6e2cux3068ux5b9fux969bux306eux7d50ux679cux306eux6bd4ux8f03-ux63cfux753bux7d50ux679c}{}
\scriptsize

\begin{center}\includegraphics[height=0.75\textheight]{QSS_Ch04_1_files/figure-beamer/unnamed-chunk-13-1} \end{center}
\normalsize
\end{frame}

\hypertarget{ux307eux3068ux3081}{%
\section{4.1.4 まとめ}\label{ux307eux3068ux3081}}

\begin{frame}[fragile]{このセクションのまとめ (1)}
\protect\hypertarget{ux3053ux306eux30bbux30afux30b7ux30e7ux30f3ux306eux307eux3068ux3081-1}{}
\begin{itemize}
\tightlist
\item
  \textbf{予測のプロセス}:
  データのクリーニング（日付処理など）、サブセットの抽出、要約統計量の計算という一連の流れ。
\item
  \textbf{プログラミングの道具}:

  \begin{itemize}
  \tightlist
  \item
    \texttt{for} ループ:
    大量（今回は51州分）の繰り返し処理を自動化する。
  \item
    \texttt{if} 文: 条件によって処理を分岐させる。
  \end{itemize}
\end{itemize}
\end{frame}

\begin{frame}[fragile]{このセクションのまとめ (2)}
\protect\hypertarget{ux3053ux306eux30bbux30afux30b7ux30e7ux30f3ux306eux307eux3068ux3081-2}{}
\begin{itemize}
\tightlist
\item
  \textbf{評価}: 実際の値と比較して、誤差（平均誤差や
  RMSE）を算出することで予測精度を客観的に評価する。
\item
  \textbf{Rの操作}: \texttt{as.Date()}, \texttt{for}, \texttt{if},
  \texttt{unique()}, \texttt{min()}, \texttt{mean()}, \texttt{text()}。
\end{itemize}
\end{frame}

\end{document}
